\begingroup
\makeatletter\tagsleft@true\let\veqno\@@leqno\makeatother
\NoSourceHead
\SourceLeaf{294}
\pdfbookmark[0]{Zur Theorie der transscendenten Functionen gehörig}{x1-transscendenten}
\begin{center}\large\bfseries ZUR THEORIE\\[.6em]DER TRANSSCENDENTEN FUNCTIONEN GEHÖRIG.\end{center}
\SourceGroup{I}
\begin{center}\large [ENTWICKELUNGEN UND UMFORMUNGEN DER UNENDLICHEN\\ REIHEN UND PRODUKTE.]\end{center}
\HistoricCitation{[Aus Handbuch 16, Bb, Den astronomischen Wissenschaften gewidmet, November 1801.]}
\NumberedItem{1}
\noindent {\small[S. 40]}\par
\par\indent Es sei
\[\sum e^{-\alpha(k+\omega)^2}=T,\]
\noindent indem für $k$ alle ganze positive und negative Zahlen gesetzt werden, und
\[T=A+2B\cos\omega P+2C\cos2\omega P+2D\cos3\omega P+\text{etc.},\]
\noindent so ist
\[\begin{gathered}
A=\int T d\omega,\quad B=\int T\cos\omega P.d\omega,\\
C=\int T\cos2\omega P.d\omega,\quad D=\int T\cos3\omega P.d\omega\ \text{u. s. w.},
\end{gathered}\]
\noindent alle Integrale von $\omega=0$ bis $\omega=1$ ausgedehnt. Es ist aber klar, dass jene Integrale zwischen diesen Grenzen mit den Integralen
\[\begin{gathered}
\int e^{-\alpha\omega\omega}d\omega,\quad\int e^{-\alpha\omega\omega}\cos\omega P.d\omega,\\
\int e^{-\alpha\omega\omega}\cos2\omega P.d\omega,\quad\int e^{-\alpha\omega\omega}\cos3\omega P.d\omega\ \text{etc.}
\end{gathered}\]
\noindent übereinkommen, wenn diese von $\omega=-\infty$ bis $\omega=+\infty$ ausgedehnt werden. Da nun allgemein
\clearpage
\ArithmeticHead{288}{L}{ANALYSIS. NACHLASS.}
\SourceLeaf{295}
\[\begin{aligned}
&e^{-\alpha\omega\omega}\cos n\omega P\\
&=\tfrac12e^{-\alpha\omega\omega+in\omega P}+\tfrac12e^{-\alpha\omega\omega-in\omega P}\\
&=\tfrac12e^{-\frac{nnPP}{4\alpha}}\left(e^{-\alpha\left(\omega-\frac{inP}{2\alpha}\right)^2}+e^{-\alpha\left(\omega+\frac{inP}{2\alpha}\right)^2}\right),
\end{aligned}\]
\noindent so folgt leicht, dass
\[\int e^{-\alpha\omega\omega}\cos n\omega P.d\omega=e^{-\frac{nn\pi\pi}{\alpha}}\surd\frac\pi\alpha,\]
\noindent folglich
\[\begin{aligned}
T=\surd\frac\pi\alpha\Bigl(&1+2e^{-\frac{\pi\pi}{\alpha}}\cos\omega P+2e^{-4\frac{\pi\pi}{\alpha}}\cos2\omega P\\
&+2e^{-9\frac{\pi\pi}{\alpha}}\cos3\omega P+\text{etc.}\Bigr).
\end{aligned}\]
\par\indent In einer andern Form so
\begin{equation*}\tag*{[1]}
\sum e^{-\alpha\,k+\omega)^2}=\surd\frac\pi\alpha.e^{-\alpha\omega\omega}.\sum e^{-\frac{\pi\pi}{\alpha}\left(k+\frac{\alpha\omega i}{\pi}\right)^2}.
\end{equation*}
\noindent Allgemeiner
\[\begin{aligned}
&\sum e^{-\alpha(k+\omega)^2}\bigl(\cos2(k+\omega)\alpha\psi+i\sin2(k+\omega)\alpha\psi\bigr)\\
&=\surd\frac\pi\alpha.\sum e^{-\frac{\pi\pi}{\alpha}\left(k-\frac{\alpha\psi}{\pi}\right)^2}\bigl(\cos2k\omega\pi-i\sin2k\omega\pi\bigr)[\text{*)}].
\end{aligned}\]
\noindent Oder, $\alpha\alpha'=\pi\pi$ und $-\alpha\psi=\omega'\pi$ gesetzt,
\[\begin{aligned}
&\sum e^{-\alpha(k+\omega)^2}\bigl(\cos(2k+2\omega)\omega'\pi-i\sin(2k+2\omega)\omega'\pi\bigr)\\
&=\surd\frac\pi\alpha.\sum e^{-\alpha'(k+\omega')^2}\bigl(\cos2k\omega\pi-i\sin2k\omega\pi\bigr),\\[1em]
&\sum e^{-\alpha\left(k+\omega-\frac14\right)^2}-\sum e^{-\alpha\left(k+\omega+\frac14\right)^2}\\
&=4\surd\frac\pi\alpha.\left(e^{-\frac{\pi\pi}{\alpha}}\sin\omega P-e^{-9\frac{\pi\pi}{\alpha}}\sin3\omega P+\text{etc.}\right)[\text{**)}].
\end{aligned}\]
\ClosingRule
\Needspace{4\baselineskip}
\par\smallskip\noindent\rule{.25\linewidth}{.4pt}\par
{\small
\par\indent [*) Links vom Gleichheitszeichen hat die Handschrift $-i\sin2(k+\omega)\alpha\psi$.]
\par\indent [** Rechts vom Gleichheitszeichen hat die Handschrift $-3e^{-9\frac{\pi\pi}{\alpha}}\sin3\omega P$.]
\par}
\clearpage
\ArithmeticHead{289}{R}{{\scriptsize ZUR THEORIE DER TRANSSCENDENTEN FUNCTIONEN. I.}}
\SourceLeaf{296}
\par\indent Man kann den Lehrsatz [1, S. 288] auch so ausdrücken:
\[\sum t e^{-\pi t^4kk-2\pi tt k\surd u-\frac12\pi u}\]
\noindent ändert den Werth nicht, wenn $t$ in $\dfrac1t$ und $u$ in $-u$ verwandelt wird.
\ClosingRule
\par\indent Das Arithmetisch-Geometrische Mittel zwischen
\[(1+2x+\ldots)^2\quad\text{und}\quad\left(2x^{\frac14}+2x^{\frac94}+\text{u. s. w.}\right)^2\]
\noindent ist $\dfrac\pi{\log\dfrac1x}\mathpunct{\raisebox{.5ex}{.}}$
\ClosingRule
\noindent {\small[S. 41]}\par
\par\indent Das Integral
\[\int\frac{d\varphi}{\surd(aa\cos\varphi^2+bb\sin\varphi^2)}\]
\noindent wird gleich dem Integrale
\[\int\frac{d\varphi'}{\surd(a'a'\cos\varphi'^2+b'b'\sin\varphi'^2)},\]
\noindent beide von $\varphi$ und $\varphi'=0$ gerechnet, wenn man
\[a'=\tfrac12(a+b),\quad b'=\surd ab\]
\noindent und
\[\frac1{\sin\varphi}=\frac{a+b}{2a}.\frac1{\sin\varphi'}+\frac{a-b}{2a}.\sin\varphi'\]
\[\begin{aligned}
\sin\varphi'&=\frac{a-\surd(aa\cos\varphi^2+bb\sin\varphi^2)}{(a-b)\sin\varphi}\\
&=\frac{(a+b)\sin\varphi}{a+\surd(aa\cos\varphi^2+bb\sin\varphi^2)}\text{\textsuperscript{*)}}
\end{aligned}\]
\noindent [setzt.]
\Needspace{7\baselineskip}
\par\smallskip\noindent\rule{.25\linewidth}{.4pt}\par
{\small
\par\indent *) [mit anderer Schrift und Tinte an der Seite]
\[\frac{a-b}{a+b}=\mu,\quad\frac{2\surd\mu'}{1+\xi'}=\mu,\quad\surd\left(1-\frac{bb}{aa}\right)=\mu,\]
\[\frac2{\mu\sin\varphi}=\sin\varphi'.\surd\mu'+\frac1{\sin\varphi'.\surd\mu'}.\]
\par}
\par\smallskip\noindent{\small X 1.}\hfill 37\par

\clearpage
\endgroup
